
\section{Deep Reinforcement Learning Applications}

Deep Reinforcement Learning (DRL) represents a paradigm shift in microgrid control, offering a promising alternative that addresses the limitations of both traditional optimization and heuristic methods. DRL agents operate structurally as model-free controllers. They learn an optimal policy through continuous interaction with the environment, gradually refining their decision-making process based on a scalar reward signal. This capability allows them to internalize the complex, non-linear dynamics of the microgrid, such as battery degradation and stochastic load patterns, without requiring an explicit mathematical model of the system physics \cite{Mocanu2018}.

Early applications of Reinforcement Learning in power systems relied on value-based methods such as Q-Learning and Deep Q-Networks (DQN). Zhang et al. \cite{Zhang2018} successfully applied DQN to reduce the energy costs of a microgrid. However, DQN is inherently limited to discrete action spaces. This necessitates the discretization of continuous variables, such as battery charging power, which leads to the "curse of dimensionality" and suboptimal, bang-bang control behavior that can accelerate equipment wear.

To overcome these limitations, recent research has pivoted toward Actor-Critic architectures capable of handling continuous action spaces. Proximal Policy Optimization (PPO) \cite{Raffin2021} has gained popularity due to its numerical stability and ease of implementation. PPO uses a clipped objective function to prevent distinct policy updates that could destabilize training, making it a robust baseline for many control tasks.

More recently, the Soft Actor-Critic (SAC) algorithm \cite{Haarnoja2018} has emerged as a state-of-the-art method for continuous control. Unlike PPO, SAC optimizes a maximum entropy objective, which encourages the agent to maximize not only the expected reward but also the entropy of the policy. This entropy term acts as a rigorous regularizer, promoting diverse exploration and preventing the agent from converging prematurely to suboptimal local minima. Zhou et al. \cite{Zhou2019} demonstrated that SAC outperforms other algorithms in multi-objective energy management tasks by maintaining a more robust policy in the face of uncertainty. This thesis builds upon these findings, proposing a specialized "Robust SAC" formulation designed specifically to mitigate the impact of forecast errors.
